% GENERATED by scripts/gen_blueprint.py — do not edit.
\chapter{The three critical families}
\label{bp:ch_families}

\section{$\mathrm{AC}_n$ is 3-critical; its basic stats}

\begin{theorem}[\code{AC\_kcritical3}]
  \label{res:AC_kcritical3}
  \uses{def:AC, def:arc, def:backedge, def:cayley, def:kcritical, def:ltp, def:omegabar, def:subdel, def:tournament}
  \rocq{Digraph.applications.k5.acn_base.AC_kcritical3}
  \rocqok

\begin{verbatim}
Theorem AC_kcritical3 : (3 <= m)%N -> kcritical 3 (AC m' : tournament).
\end{verbatim}
\href{https://github.com/LLM4Rocq/digraph-theory/blob/main/theories/applications/k5/acn_base.v\#L431}{Source: \code{theories/applications/k5/acn\_base.v:431}}


\paragraph{Decoded.} $\mathrm{AC}_n$ is $3$-$\bar\omega$-critical (\ref{def:kcritical}),
has exactly $n$ vertices, and is vertex-transitive (\ref{def:vt}).

\end{theorem}

\section{$\bar\omega(\mathrm{AC}_n) = 3$}

\begin{theorem}[\code{omegabar\_AC}]
  \label{res:omegabar_AC}
  \uses{def:AC, def:arc, def:backedge, def:cayley, def:ltp, def:omegabar, def:tournament}
  \rocq{Digraph.applications.k5.acn_base.omegabar_AC}
  \rocqok

\begin{verbatim}
Theorem omegabar_AC : (3 <= m)%N -> ω̄((AC m' : tournament)) = 3.
\end{verbatim}
\href{https://github.com/LLM4Rocq/digraph-theory/blob/main/theories/applications/k5/acn_base.v\#L374}{Source: \code{theories/applications/k5/acn\_base.v:374}}


\paragraph{Decoded.} For every odd $n = 2m+1$ with $m \ge 3$, the tournament clique number
of $\mathrm{AC}_n$ is exactly $3$. (\code{(AC m' : tournament)}
is $\mathrm{AC}_n$ viewed as a tournament; the hypothesis
\code{3 <= m} is $m \ge 3$.)

\end{theorem}

\section{$\bar\omega(\mathrm{AC}_n - v) = 2$ for every $v$}

\begin{theorem}[\code{omegabar\_AC\_del}]
  \label{res:omegabar_AC_del}
  \uses{def:AC, def:arc, def:backedge, def:cayley, def:ltp, def:omegabar, def:subdel, def:tournament}
  \rocq{Digraph.applications.k5.acn_base.omegabar_AC_del}
  \rocqok

\begin{verbatim}
Theorem omegabar_AC_del (v : ACm) :
  (3 <= m)%N -> ω̄(del_tournament v) = 2.
\end{verbatim}
\href{https://github.com/LLM4Rocq/digraph-theory/blob/main/theories/applications/k5/acn_base.v\#L423}{Source: \code{theories/applications/k5/acn\_base.v:423}}


\paragraph{Decoded.} Deleting any single vertex $v$ from $\mathrm{AC}_n$ drops the
invariant to $2$.

\end{theorem}

\begin{theorem}[\code{card\_AC}]
  \label{res:card_AC}
  \uses{def:AC, def:cayley, def:tournament}
  \rocq{Digraph.constructions.circulant.card_AC}
  \rocqok

\begin{verbatim}
Lemma card_AC : #|AC| = n.
\end{verbatim}
\href{https://github.com/LLM4Rocq/digraph-theory/blob/main/theories/constructions/circulant.v\#L129}{Source: \code{theories/constructions/circulant.v:129}}


\end{theorem}

\begin{theorem}[\code{AC\_vertex\_transitive}]
  \label{res:AC_vertex_transitive}
  \uses{def:AC, def:arc, def:aut, def:cayley, def:tournament, def:vt}
  \rocq{Digraph.constructions.circulant.AC_vertex_transitive}
  \rocqok

\begin{verbatim}
Lemma AC_vertex_transitive : vertex_transitiveb AC.
\end{verbatim}
\href{https://github.com/LLM4Rocq/digraph-theory/blob/main/theories/constructions/circulant.v\#L134}{Source: \code{theories/constructions/circulant.v:134}}


\end{theorem}

\section{$\mathrm{AC}_n[C_3]$ is 4-critical, on $3n$ vertices}

\begin{theorem}[\code{T4\_kcritical4}]
  \label{res:T4_kcritical4}
  \uses{def:AC, def:C3, def:T4family, def:arc, def:backedge, def:cayley, def:kcritical, def:lexprod, def:ltp, def:omegabar, def:subdel, def:tournament}
  \rocq{Digraph.applications.k4.k4_main.T4_kcritical4}
  \rocqok

\begin{verbatim}
Theorem T4_kcritical4 : kcritical 4 T4.
\end{verbatim}
\href{https://github.com/LLM4Rocq/digraph-theory/blob/main/theories/applications/k4/k4_main.v\#L59}{Source: \code{theories/applications/k4/k4\_main.v:59}}


\paragraph{Decoded.} The two previous values say precisely that $\mathrm{AC}_n[C_3]$ is
$4$-$\bar\omega$-critical; its order is $n \cdot 3$.

\end{theorem}

\section{$\bar\omega(\mathrm{AC}_n[C_3]) = 4$}

\begin{theorem}[\code{omegabar\_T4}]
  \label{res:omegabar_T4}
  \uses{def:AC, def:C3, def:T4family, def:arc, def:backedge, def:cayley, def:lexprod, def:ltp, def:omegabar, def:tournament}
  \rocq{Digraph.applications.k4.k4_value.omegabar_T4}
  \rocqok

\begin{verbatim}
Theorem omegabar_T4 : ω̄(T4) = 4.
\end{verbatim}
\href{https://github.com/LLM4Rocq/digraph-theory/blob/main/theories/applications/k4/k4_value.v\#L511}{Source: \code{theories/applications/k4/k4\_value.v:511}}


\paragraph{Decoded.} Here \code{T4} abbreviates the substitution
$\mathrm{AC}_n[C_3]$ (\ref{def:T4family}). For every odd
$n = 2m{+}1 \ge 7$, its tournament clique number is exactly $4$.

\end{theorem}

\section{$\bar\omega(\mathrm{AC}_n[C_3] - v) = 3$ for every $v$}

\begin{theorem}[\code{omegabar\_T4\_del}]
  \label{res:omegabar_T4_del}
  \uses{def:AC, def:C3, def:T4family, def:arc, def:backedge, def:cayley, def:lexprod, def:ltp, def:omegabar, def:subdel, def:tournament}
  \rocq{Digraph.applications.k4.k4_main.omegabar_T4_del}
  \rocqok

\begin{verbatim}
Theorem omegabar_T4_del (v : T4) : ω̄(del_tournament v) = 3.
\end{verbatim}
\href{https://github.com/LLM4Rocq/digraph-theory/blob/main/theories/applications/k4/k4_main.v\#L51}{Source: \code{theories/applications/k4/k4\_main.v:51}}


\paragraph{Decoded.} Deleting any vertex of $\mathrm{AC}_n[C_3]$ drops the invariant to
$3$ (by vertex-transitivity all deletions are isomorphic, so this is
one computation plus symmetry --- see \ref{res:vt_kcritical}).

\end{theorem}

\begin{theorem}[\code{card\_T4}]
  \label{res:card_T4}
  \uses{def:AC, def:C3, def:T4family, def:cayley, def:lexprod, def:tournament}
  \rocq{Digraph.applications.k4.k4_main.card_T4}
  \rocqok

\begin{verbatim}
Lemma card_T4 : #|T4| = (n * 3)%N.
\end{verbatim}
\href{https://github.com/LLM4Rocq/digraph-theory/blob/main/theories/applications/k4/k4_main.v\#L67}{Source: \code{theories/applications/k4/k4\_main.v:67}}


\end{theorem}

\section{$\bar\omega(\mathrm{AC}_n[\mathrm{AC}_n] - v) = 4$; 5-criticality; order $n^2$}

\begin{theorem}[\code{T5\_kcritical5}]
  \label{res:T5_kcritical5}
  \uses{def:AC, def:T5family, def:arc, def:backedge, def:cayley, def:kcritical, def:lexprod, def:ltp, def:omegabar, def:subdel, def:tournament}
  \rocq{Digraph.applications.k5.main.T5_kcritical5}
  \rocqok

\begin{verbatim}
Theorem T5_kcritical5 : kcritical 5 T5.
\end{verbatim}
\href{https://github.com/LLM4Rocq/digraph-theory/blob/main/theories/applications/k5/main.v\#L71}{Source: \code{theories/applications/k5/main.v:71}}


\end{theorem}

\section{$\bar\omega(\mathrm{AC}_n[\mathrm{AC}_n]) = 5$}

\begin{theorem}[\code{omegabar\_T5}]
  \label{res:omegabar_T5}
  \uses{def:AC, def:T5family, def:arc, def:backedge, def:cayley, def:lexprod, def:ltp, def:omegabar, def:tournament}
  \rocq{Digraph.applications.k5.main.omegabar_T5}
  \rocqok

\begin{verbatim}
Theorem omegabar_T5 : ω̄(T5) = 5.
\end{verbatim}
\href{https://github.com/LLM4Rocq/digraph-theory/blob/main/theories/applications/k5/main.v\#L54}{Source: \code{theories/applications/k5/main.v:54}}


\paragraph{Decoded.} \code{T5} abbreviates $\mathrm{AC}_n[\mathrm{AC}_n]$
(\ref{def:T5family}); its tournament clique number is exactly $5$,
for every odd $n = 2m{+}1 \ge 7$.

\end{theorem}

\begin{theorem}[\code{omegabar\_T5\_del}]
  \label{res:omegabar_T5_del}
  \uses{def:AC, def:T5family, def:arc, def:backedge, def:cayley, def:lexprod, def:ltp, def:omegabar, def:subdel, def:tournament}
  \rocq{Digraph.applications.k5.main.omegabar_T5_del}
  \rocqok

\begin{verbatim}
Theorem omegabar_T5_del (v : T5) : ω̄(del_tournament v) = 4.
\end{verbatim}
\href{https://github.com/LLM4Rocq/digraph-theory/blob/main/theories/applications/k5/main.v\#L62}{Source: \code{theories/applications/k5/main.v:62}}


\paragraph{Decoded.} Every single-vertex deletion has $\bar\omega = 4$; hence
$\mathrm{AC}_n[\mathrm{AC}_n]$ is $5$-$\bar\omega$-critical, on
$n \cdot n$ vertices.

\end{theorem}

\begin{theorem}[\code{card\_T5}]
  \label{res:card_T5}
  \uses{def:AC, def:T5family, def:cayley, def:lexprod, def:tournament}
  \rocq{Digraph.applications.k5.main.card_T5}
  \rocqok

\begin{verbatim}
Lemma card_T5 : #|T5| = (n * n)%N.
\end{verbatim}
\href{https://github.com/LLM4Rocq/digraph-theory/blob/main/theories/applications/k5/main.v\#L79}{Source: \code{theories/applications/k5/main.v:79}}


\end{theorem}

